\documentclass[11pt]{article}
\usepackage{ap-physc}
\course{AP Physics C: Electricity and Magnetism}
\headcourse{AP Physics C: E\&M}
\headpart{Unit 9 \textperiodcentered\ Question Set}
\title{Unit 9: Electric Potential: Question Set}
\author{Paul Wang}
\date{November 2026}
\docpart{Unit 9 \textperiodcentered\ Question Set}

\begin{document}
\maketitle

\begin{abstract}
Exam-style practice on Unit 9 (topics 9.1--9.3), which is 10--20\% of the multiple-choice section. Section~1--3 hold 12 four-choice multiple-choice questions grouped by topic; most need a short calculation or a line of reasoning, as on the exam. The last page before the key is a full Qualitative/Quantitative Translation free-response question (8 points). All items are original and written for practice. Use the AP equation sheet and a calculator. Allow about 24 minutes for the multiple choice (the exam's pace is about 2 minutes per question) and 24 minutes for the free response. The answer key explains every item and names its topic; a scoring guide for the free response follows it.
\end{abstract}

\scoringline{MCQ: \blank\ / \qtotal \qquad FRQ: \blank\ / 8 \qquad Target: 8+ MCQ \qquad Pace: 2 min per MCQ}

\section{Potential energy (9.1)}

\begin{q}{2}{9.1 Energy of a system}{There are three pairs, each separated by $a$: $U = 3 \times kq^2/a$.}
Three charges $+q$ are placed at the corners of an equilateral triangle of side $a$. What is the electric potential energy of the system?
\chh{$kq^2/a$}{$3kq^2/a$}{$6kq^2/a$}{$3kq^2/a^2$}
\end{q}
\begin{q}{3}{9.1 Uniform field}{$\Delta U = -qEd = -(1.6\times10^{-19})(200)(0.050) = -1.6\times10^{-18}$ J: the field does positive work, so $U$ falls.}
A proton moves 0.050 m in the direction of a uniform electric field of 200 N/C. What is the change in the electric potential energy of the proton--field system?
\chh{$+1.6\times10^{-18}$ J}{$0$}{$-1.6\times10^{-18}$ J}{$-3.2\times10^{-17}$ J}
\end{q}

\section{Electric potential (9.2)}

\begin{q}{3}{9.2 Potential of a point charge}{$V = kq/r = (9.0\times10^9)(2.0\times10^{-9})/0.30 = 60$ V.}
What is the electric potential 0.30 m from a $+2.0$ nC point charge? ($k = 9.0\times10^9\ \mathrm{N\,m^2/C^2}$)
\chh{$5.4$ V}{$18$ V}{$60$ V}{$200$ V}
\end{q}
\begin{q}{4}{9.2 Dipole midpoint}{The potentials $+kq/(d/2)$ and $-kq/(d/2)$ cancel. The fields both point toward $-q$ and add to $8kq/d^2$.}
Charges $+q$ and $-q$ are a distance $d$ apart. At the midpoint between them, which is true?
\chh{$V = 0$ and $E = 0$}{$V \ne 0$ and $E = 0$}{$V \ne 0$ and $E \ne 0$}{$V = 0$ and $E \ne 0$}
\end{q}
\begin{q}{2}{9.2 Potential difference from a field}{$\Delta V = -\int_0^3 4x^2\,dx = -\tfrac43(27) = -36$ V. The field points in $+x$, so the potential falls along $+x$.}
The electric field along the $x$-axis is $\vec E = 4x^2\,\hat\imath$ (N/C, $x$ in m). What is $V(3.0\ \mathrm{m}) - V(0)$?
\chh{$-108$ V}{$-36$ V}{$-18$ V}{$+36$ V}
\end{q}
\begin{q}{2}{9.2 Field from potential}{$E_x = -dV/dx = -(15x^2 - 2) = -13$ V/m at $x = 1$.}
The potential along the $x$-axis is $V(x) = 5x^3 - 2x$ (V, $x$ in m). What is $E_x$ at $x = 1.0$ m?
\chh{$-15$ V/m}{$-13$ V/m}{$+3.0$ V/m}{$+13$ V/m}
\end{q}
\begin{q}{1}{9.2 Potential of a ring}{Every bit of charge is the same distance $R$ from the center, and potential is a scalar: $V = \int k\,dq/R = kQ/R$, however the charge is spread.}
A ring of radius $R$ carries total charge $Q$ spread unevenly around it. What is the electric potential at its center?
\ch{$kQ/R$}{$0$, because the charge is spread around a circle}{$kQ/R^2$}{It cannot be found without knowing how the charge is distributed.}
\end{q}
\begin{q}{4}{9.2 Equipotentials}{$E = |\Delta V|/d = 10/0.020 = 500$ V/m, pointing toward lower potential, west.}
Equipotential lines 10 V apart are parallel and 2.0 cm apart, with the potential increasing toward the east. What is the electric field?
\chh{$500$ V/m, pointing east}{$5.0$ V/m, pointing west}{$20$ V/m, pointing east}{$500$ V/m, pointing west}
\end{q}
\begin{q}{1}{9.2 Zero potential}{The field depends on how $V$ changes ($E = -dV/dx$), not on its value. The midpoint of a dipole has $V = 0$ and $E \ne 0$; a point far from all charges has both zero.}
At a certain point the electric potential is zero. What must be true of the electric field there?
\ch{Nothing; the field can be zero or nonzero.}{The field must be zero.}{The field must be nonzero.}{The field must point toward the nearest negative charge.}
\end{q}

\section{Conservation of energy (9.3)}

\begin{q}{3}{9.3 Accelerating an electron}{$\tfrac12 mv^2 = e\Delta V$, so $v = \sqrt{2e\Delta V/m} = \sqrt{1.76\times10^{14}} = 1.3\times10^7$ m/s.}
An electron starts from rest and is accelerated through a potential difference of 500 V. What is its final speed? ($e = 1.6\times10^{-19}$ C, $m_e = 9.1\times10^{-31}$ kg)
\chh{$3.1\times10^{5}$ m/s}{$9.4\times10^{6}$ m/s}{$1.3\times10^{7}$ m/s}{$8.8\times10^{13}$ m/s}
\end{q}
\begin{q}{4}{9.3 Closest approach}{All kinetic energy becomes potential energy at the turning point: $\tfrac12 mv_0^2 = kqQ/r_{min}$.}
A particle of mass $m$ and charge $+q$ is fired with speed $v_0$ straight at a fixed charge $+Q$ from far away. How close does it get?
\ch{$\dfrac{kqQ}{mv_0^2}$}{$\sqrt{\dfrac{2kqQ}{mv_0^2}}$}{$\dfrac{mv_0^2}{2kqQ}$}{$\dfrac{2kqQ}{mv_0^2}$}
\end{q}
\begin{q}{1}{9.3 Electron moving between potentials}{$\Delta U = q\Delta V = (-e)(+200\ \mathrm{V}) = -200$ eV, so $\Delta K = +200$ eV. Electrons gain speed moving toward higher potential.}
An electron moves, with only electric forces acting, from a point at $V = 100$ V to a point at $V = 300$ V. Its kinetic energy
\chh{increases by 200 eV}{decreases by 200 eV}{increases by 300 eV}{decreases by 400 eV}
\end{q}

\section{Free response}

\begin{frq}{Question 1 \textperiodcentered\ Qualitative/Quantitative Translation \textperiodcentered\ 8 points \textperiodcentered\ 24 minutes}
Two identical positive charges $+Q$ are fixed at $x = -a$ and $x = +a$. A small particle of mass $m$ and charge $+q$ can move freely along the $x$-axis between them. Ignore gravity.

\begin{enumerate}[label=(\alph*),leftmargin=2em,itemsep=4pt]
\item The particle is released from rest at $x = a/2$. Without equations, describe its subsequent motion and explain where its speed is greatest. \pts{2}
\item Write an expression for the electric potential $V(x)$ for $-a < x < a$. \pts{2}
\item Derive an expression for the particle's speed as it passes $x = 0$, in terms of $k$, $Q$, $q$, $m$, and $a$. \pts{3}
\item Explain how the expression for $V(x)$ in part (b) supports your claim in part (a) about where the speed is greatest. \pts{1}
\end{enumerate}
\end{frq}

\makekey

\subsection*{Free-response scoring guide (8 points)}
\small
Award each point only when the answer does what the model does. Points are independent unless a part says otherwise; a correct method with an answer carried forward from an earlier error earns the later points.
\begin{scoring}
(a) (2) & 1 pt: the particle is pushed back toward the center, passes through $x = 0$, slows, stops at $x = -a/2$, and returns: it oscillates between $\pm a/2$. 1 pt: its speed is greatest at $x = 0$, where the net force is zero (forces from the two charges balance) and the potential, hence potential energy, is lowest along the segment. & Saying the particle stops at the center because the net force is zero there. \\
(b) (2) & 1 pt: superposition of two point-charge potentials. 1 pt: $V(x) = \dfrac{kQ}{a + x} + \dfrac{kQ}{a - x}$ (equivalently $\dfrac{2kQa}{a^2 - x^2}$). & Subtracting the two terms; both charges are positive. \\
(c) (3) & 1 pt: $V(a/2) = kQ\left(\dfrac{2}{3a} + \dfrac{2}{a}\right) = \dfrac{8kQ}{3a}$ and $V(0) = \dfrac{2kQ}{a}$. 1 pt: energy conservation, $q[V(a/2) - V(0)] = \tfrac12 mv^2$, with $V(a/2) - V(0) = \dfrac{2kQ}{3a}$. 1 pt: $v = \sqrt{\dfrac{4kQq}{3ma}}$. & Using the force at $x = a/2$ with constant-acceleration kinematics. \\
(d) (1) & 1 pt: $V(x) = 2kQa/(a^2 - x^2)$ has its minimum at $x = 0$ (its slope, $-E_x$, is zero there and it rises toward both charges), so a positive charge has its lowest potential energy and greatest kinetic energy at $x = 0$. & Pointing to part (c) without referring to the shape of $V(x)$. \\
\end{scoring}

\end{document}
